\documentclass[a4paper, serif, 10pt]{moderncv}

%% moderncv
% color and style
\moderncvcolor{black}
\moderncvstyle{banking}

%% page layout/margins
\usepackage[top=2.5cm, bottom=2.5cm, left=1.5cm, right=1.5cm]{geometry}

\nopagenumbers{}

%% Babel config for Brazilian Portuguese language
% ref:
% - https://latex3.github.io/babel/guides/locale-portuguese.html
\usepackage[brazilian]{babel}

%% fontspec
\usepackage{fontspec}

\IfFontExistsTF{Linux Libertine}{
  \setmainfont[Ligatures={TeX, Common}, Numbers=OldStyle]{Linux Libertine}
}{}
\IfFontExistsTF{Linux Libertine O}{
  \setmainfont[Ligatures={TeX, Common}, Numbers=OldStyle]{Linux Libertine O}
}{}

%% CTeX
% CJK support, used to show CJK characters in the resume
%
% - heading=false: disable ctex heading
% - punct=kaiming: use kaiming punctuations styles for CJK
% - scheme=plain: use plain scheme, do not override \normalsize font size
% - space=auto: space settings for CJK characters
%
% CJK fonts are resolved per locale language, e.g. Noto Serif/Sans CJK JP
% for Japanese resumes, with fallback to the next region if not installed
%
% ref:
% - http://ctan.mirrorcatalogs.com/language/chinese/ctex/ctex.pdf
\usepackage[UTF8, heading=false, punct=kaiming, scheme=plain, space=auto]{ctex}

\IfFontExistsTF{Noto Serif CJK SC}{
  \setCJKmainfont{Noto Serif CJK SC}
}{}
\IfFontExistsTF{Noto Sans CJK SC}{
  \setCJKsansfont{Noto Sans CJK SC}
}{}

%% line spacing
\usepackage{setspace}
\setstretch{1.125}

%% URL styling - use normal text instead of monospace
\urlstyle{same}

%% Patch \httplink and \httpslink to handle full URLs
% The original moderncv macros always prepend http:// or https:// to URLs.
% This causes issues when the URL already has a protocol (e.g., https://example.com)
% resulting in malformed URLs like https://https://example.com.
% This patch checks if the URL already contains "://" and uses it directly if so.
% Note: We use \str_set:Nx (x-expansion) to expand macros like \@homepage first.
\ExplSyntaxOn
\str_new:N \g_yamlresume_url_str

\RenewDocumentCommand{\httpslink}{O{}m}{
  \str_set:Nx \g_yamlresume_url_str {#2}
  \str_if_in:NnTF \g_yamlresume_url_str {://}
    { \url{#2} }
    { \url{https://#2} }
}

\RenewDocumentCommand{\httplink}{O{}m}{
  \str_set:Nx \g_yamlresume_url_str {#2}
  \str_if_in:NnTF \g_yamlresume_url_str {://}
    { \url{#2} }
    { \url{http://#2} }
}
\ExplSyntaxOff

\name{Mariana Almeida Costa}{}
\title{Enfermeira Registrada | Terapia Intensiva e Segurança do Paciente}
\phone[mobile]{+55 11 98765-4321}
\email{mariana.costa@enfermagem.com.br}
\homepage{https://marianacosta-enfermagem.com.br}

\address{Rua Cardoso de Almeida, 1234, Apto 56, São Paulo, SP, Brasil, 01215-000}{}{}

\extrainfo{{\small \faLinkedin}\ \href{https://www.linkedin.com/in/mariana-almeida-costa-rn}{@mariana-almeida-costa-rn} {} {} {} • {} {} {} 
{\small \faInstagram}\ \href{https://www.instagram.com/enf.marianacosta}{@enf.marianacosta}}

\begin{document}

\maketitle

\section{Informações Básicas}

\cvline{}{\begin{itemize}
\item Enfermeira registrada com mais de 6 anos de experiência em assistência hospitalar, atuando em unidades de terapia intensiva e clínica médica.
\item Sólida formação em enfermagem clínica, com foco em segurança do paciente, prevenção de infecções e aplicação da Sistematização da Assistência de Enfermagem (SAE).
\item Experiência na liderança de equipes de técnicos e auxiliares de enfermagem, elaboração de escalas, educação continuada e monitoramento de indicadores assistenciais.
\item Comunicação clara e empática no cuidado ao paciente e à família, com atuação colaborativa em equipes multiprofissionais.
\item Registro profissional ativo no COREN-SP, com prática alinhada às boas práticas de enfermagem e às normas de biossegurança.
\end{itemize}}
\section{Formação}

\cventry{fev. de 2014–dez. de 2017}
        {Bacharelado, Enfermagem, Nota: 8.7/10}
        {Universidade de São Paulo}
        {\url{https://www.usp.br}}
        {}
        {\begin{itemize}
\item Formação generalista com ênfase em enfermagem clínica, saúde coletiva e gestão do cuidado.
\item Desenvolvimento de competências para avaliação integral do paciente, raciocínio clínico e tomada de decisão baseada em evidências.
\item Participação em projetos de extensão e atividades práticas em hospitais universitários.
\end{itemize}
\textbf{Cursos}: Anatomia Humana, Fisiologia, Farmacologia, Semiologia e Semiotécnica, Enfermagem em Terapia Intensiva, Saúde Coletiva, Gestão em Enfermagem, Bioética}
\section{Experiência Profissional}

\cventry{jan. de 2022–Atual}
        {Enfermeira Assistencial}
        {Hospital Israelita Albert Einstein}
        {\url{https://www.einstein.br}}
        {}
        {\begin{itemize}
\item Presto assistência de enfermagem a pacientes adultos em unidade de terapia intensiva, realizando avaliação clínica, administração de medicamentos, curativos complexos e monitorização hemodinâmica.
\item Aplico protocolos institucionais de segurança do paciente, prevenção de infecção relacionada à assistência à saúde e notificação de eventos adversos.
\item Lidero equipe de técnicos de enfermagem durante o turno, organizando escalas, distribuindo atividades e apoiando a tomada de decisão clínica.
\item Participo de rondas multiprofissionais, discussões de casos e treinamentos de educação continuada voltados à qualidade do cuidado.
\item Utilizo prontuário eletrônico e sistemas assistenciais para registro, auditoria e acompanhamento de indicadores.
\end{itemize}
\textbf{Palavras-chave}: Terapia Intensiva, Segurança do Paciente, SAE, Liderança de Equipe, Prontuário Eletrônico}

\cventry{mar. de 2018–dez. de 2021}
        {Enfermeira Assistencial}
        {Hospital das Clínicas da FMUSP}
        {\url{https://www.hc.fm.usp.br}}
        {}
        {\begin{itemize}
\item Prestei cuidados de enfermagem a pacientes adultos e idosos em unidades de clínica médica e cuidados paliativos, com foco em conforto, dignidade e continuidade do cuidado.
\item Realizei triagem, acolhimento, avaliação inicial, administração de medicamentos e orientação de alta hospitalar.
\item Desenvolvi ações de educação em saúde para pacientes e familiares sobre prevenção de complicações, adesão ao tratamento e autocuidado.
\item Colaborei com equipe multiprofissional na elaboração de planos de cuidado individualizados e na comunicação de intercorrências.
\item Apoiei a supervisão de estagiários e novos colaboradores, contribuindo para a integração e capacitação da equipe.
\end{itemize}
\textbf{Palavras-chave}: Clínica Médica, Cuidados Paliativos, Educação em Saúde, Triagem}
\section{Idiomas}

\cvline{Português}{Nativo ou Bilíngue}
\cvline{Inglês}{Proficiência Profissional Fluente \hfill \textbf{Palavras-chave}: TOEFL iBT 105, Leitura de artigos científicos}
\cvline{Espanhol}{Proficiência Profissional Limitada \hfill \textbf{Palavras-chave}: DELE B1}
\section{Competências}

\cvline{Enfermagem Clínica}{Especialista \hfill \textbf{Palavras-chave}: Avaliação de Enfermagem, Administração de Medicamentos, Curativos Complexos, Cateterismo Vesical, Coleta de Exames}
\cvline{Terapia Intensiva}{Avançado \hfill \textbf{Palavras-chave}: Ventilação Mecânica, Monitorização Hemodinâmica, Gasometria Arterial, Protocolos de Sepse, Suporte Avançado de Vida}
\cvline{Segurança do Paciente}{Avançado \hfill \textbf{Palavras-chave}: Prevenção de Infecção, Higienização das Mãos, Notificação de Eventos, Auditoria de Prontuários}
\cvline{Gestão e Liderança}{Avançado \hfill \textbf{Palavras-chave}: Escalas de Trabalho, Educação Continuada, Indicadores Assistenciais, Gestão de Conflitos}
\cvline{Sistemas de Saúde}{Intermediário \hfill \textbf{Palavras-chave}: Prontuário Eletrônico, TASY, MV Soul}
\section{Prêmios}

\cventry{nov. de 2023}
        {Hospital Israelita Albert Einstein}
        {Prêmio de Excelência em Cuidado ao Paciente}
        {}
        {}
        {Reconhecimento por desempenho consistente na assistência de enfermagem, comunicação empática com pacientes e familiares e compromisso com a segurança do paciente.}

\cventry{dez. de 2017}
        {Escola de Enfermagem da USP}
        {Destaque Acadêmico em Enfermagem}
        {}
        {}
        {Concedido a estudantes com alto desempenho acadêmico e envolvimento em atividades de extensão universitária.}
\section{Certificados}

\cventry{fev. de 2024}
        {American Heart Association}
        {Certificação em Suporte Básico de Vida (BLS)}
        {\url{https://www.heart.org/en/courses/bls}}
        {}
        {}

\cventry{dez. de 2021}
        {Sociedade Brasileira de Terapia Intensiva}
        {Especialização em Enfermagem em Terapia Intensiva}
        {\url{https://www.sbti.org.br}}
        {}
        {}

\cventry{ago. de 2020}
        {Associação Brasileira de Enfermagem}
        {Curso de Interpretação de Gasometria Arterial}
        {\url{https://www.abennacional.org.br}}
        {}
        {}
\section{Publicações}

\cventry{jun. de 2021}
        {Cuidados de Enfermagem na Prevenção de Lesões por Pressão em Pacientes Críticos}
        {Revista Brasileira de Enfermagem}
        {\url{https://www.scielo.br/j/reben/}}
        {}
        {\begin{itemize}
\item Revisão integrativa sobre estratégias de prevenção de lesões por pressão em unidades de terapia intensiva.
\item Discussão de protocolos assistenciais, avaliação de risco e indicadores de qualidade do cuidado.
\end{itemize}}
\section{Projetos}

\cventry{mar. de 2022–dez. de 2023}
        {Projeto multiprofissional para reduzir eventos adversos e fortalecer a cultura de segurança em unidades de internação.}
        {Programa de Educação Continuada em Segurança do Paciente}
        {\url{https://www.einstein.br/ensino-e-pesquisa}}
        {}
        {\begin{itemize}
\item Desenvolvi materiais educativos e treinamentos sobre higienização das mãos, identificação segura do paciente e prevenção de quedas.
\item Acompanhei indicadores assistenciais e apresentei resultados em reuniões de qualidade.
\item Contribuí para a integração de novos colaboradores e para a padronização de condutas de enfermagem.
\end{itemize}
\textbf{Palavras-chave}: Segurança do Paciente, Educação Continuada, Indicadores de Qualidade}
\section{Interesses}

\cvline{Corrida de Rua}{Meia maratona, Treino funcional}
\cvline{Leitura}{Literatura brasileira, Saúde pública}
\cvline{Voluntariado}{Ações comunitárias, Educação em saúde}
\section{Voluntariado}

\cventry{mai. de 2019–nov. de 2023}
        {Enfermeira Voluntária}
        {Associação Brasileira de Enfermagem}
        {\url{https://www.abennacional.org.br}}
        {}
        {\begin{itemize}
\item Participei de ações comunitárias de promoção da saúde, prevenção de doenças e educação em saúde.
\item Apoiei campanhas de vacinação, orientação sobre higiene e cuidados básicos com a saúde da família.
\item Colaborei na organização de eventos científicos e atividades de atualização profissional.
\end{itemize}}
    
\end{document}